\begin{proof}[Proof of \cref{obj:lem:localized-clipped-drift-bound}]
Fix the evaluation fold and abbreviate the foldwise estimators by \((\widehat\mu_0,\widehat\mu_1,\widehat e)\).  Write
\[
\bar e_q(x)=\min\{1-q,\max\{q,\widehat e(x)\}\},
\qquad
\Delta_a(x)=\widehat\mu_a(x)-\mu_a(x),
\qquad
\Delta_e(x)=\widehat e(x)-e_P(x),
\]
and, for a policy \(\pi\),
\[
T_\pi(q)=D_\pi\cap\{x:p_P(x)\le q\},
\qquad
P\!\left[(\pi-\pi^\star_P)b_q\right]
=\int_{\mathcal X}\left\{\mathbf 1\{\pi(x)=1\}-\mathbf 1\{\pi^\star_P(x)=1\}\right\}b_q(x)\,dP_X(x).
\]
By the exact drift identity of \cref{obj:lem:clip-bias}, the conditional-mean drift defined in the statement has the closed form
\[
b_q(x)=\{\bar e_q(x)-e_P(x)\}\left\{\frac{\Delta_1(x)}{\bar e_q(x)}+\frac{\Delta_0(x)}{1-\bar e_q(x)}\right\}.
\]
% lean: clip_bias

\emph{Step 1: a pointwise majorization.}  We bound the two factors of \(b_q\) separately.

First, the clipping map \(z\mapsto\min\{1-q,\max\{q,z\}\}\) is \(1\)-Lipschitz, being a composition of the \(1\)-Lipschitz maps \(z\mapsto\max\{q,z\}\) and \(z\mapsto\min\{1-q,z\}\); applying it to \(\widehat e(x)\) and to \(e_P(x)\) therefore gives
\[
\left|\bar e_q(x)-\min\{1-q,\max\{q,e_P(x)\}\}\right|\le|\Delta_e(x)| .
\]
% lean: clippedPropensity_lipschitz
Second, clipping the true propensity moves it only inside the weak-overlap region: if \(p_P(x)>q\), then \(q<e_P(x)<1-q\) and the clip is inactive, so the displacement is zero; and in all cases the displacement is at most \(q\), because \(e_P(x)\in[0,1]\) and the clipped value lies in \([q,1-q]\).  Hence
\[
\left|\min\{1-q,\max\{q,e_P(x)\}\}-e_P(x)\right|\le q\,\mathbf 1\{p_P(x)\le q\},
\]
and by the triangle inequality
\[
\left|\bar e_q(x)-e_P(x)\right|\le|\Delta_e(x)|+q\,\mathbf 1\{p_P(x)\le q\} .
\]
% lean: clippedPropensity_error_le_error_plus_overlap_indicator
Third, since \(q\le\bar e_q(x)\le1-q\), both denominators are at least \(q>0\), so
\[
\left|\frac{\Delta_1(x)}{\bar e_q(x)}+\frac{\Delta_0(x)}{1-\bar e_q(x)}\right|
\le\frac{|\Delta_1(x)|+|\Delta_0(x)|}{q}.
\]
Multiplying the last two displays,
\[
|b_q(x)|
\le
\left\{|\Delta_e(x)|+q\,\mathbf 1\{p_P(x)\le q\}\right\}\cdot\frac{|\Delta_1(x)|+|\Delta_0(x)|}{q}.
\]
Finally, the policy factor satisfies \(|\mathbf 1\{\pi(x)=1\}-\mathbf 1\{\pi^\star_P(x)=1\}|\le1\), and it vanishes off \(D_\pi\); consequently
\[
\left|\mathbf 1\{\pi(x)=1\}-\mathbf 1\{\pi^\star_P(x)=1\}\right|\cdot\mathbf 1\{p_P(x)\le q\}
\le
\mathbf 1\{x\in T_\pi(q)\} .
\]
% lean: policy_overlap_indicator_mul_le_inter_indicator
Combining the last two displays and cancelling the factor \(q\) against \(q^{-1}\) in the indicator term gives the pointwise majorization
\[
\left|\left\{\mathbf 1\{\pi(x)=1\}-\mathbf 1\{\pi^\star_P(x)=1\}\right\}b_q(x)\right|
\le
\frac{|\Delta_e(x)||\Delta_1(x)|+|\Delta_e(x)||\Delta_0(x)|}{q}
+\mathbf 1\{x\in T_\pi(q)\}\left\{|\Delta_1(x)|+|\Delta_0(x)|\right\}.
\]
% lean: policy_clipBias_abs_pointwise_le

\emph{Step 2: integrate.}  All four majorant terms are \(P_X\)-integrable, being products of \(L_2(P_X)\) functions (the assumed \(L_2\) memberships of the three nuisance errors, and the indicator, which is bounded).  Integrating and applying Cauchy--Schwarz with the assumed \(L_2(P_X)\) rates gives, for \(a=0,1\),
\[
\int |\Delta_e|\,|\Delta_a|\,dP_X
   \le \left(\int\Delta_e^2\,dP_X\right)^{1/2}\left(\int\Delta_a^2\,dP_X\right)^{1/2}
   \le r_{e,n}r_{\mu,n},
\]
% lean: l2_abs_product_integral_le
and, again by Cauchy--Schwarz with the indicator in the first slot,
\[
\int \mathbf 1_{T_\pi(q)}\,|\Delta_a|\,dP_X
\le
P_X\!\left(T_\pi(q)\right)^{1/2}r_{\mu,n}.
\]
% lean: l2_set_abs_integral_le
Summing the four contributions yields the basic drift bound
\[
\left|P\!\left[(\pi-\pi^\star_P)b_q\right]\right|
\le
2\,\frac{r_{\mu,n}r_{e,n}}{q}
+
2\,r_{\mu,n}\,P_X\!\left(T_\pi(q)\right)^{1/2}.
\]
% lean: clipBias_drift_l2_mass_bound

\emph{Step 3: the case \(\gamma>0\).}  Assume \(\gamma>0\), \(0<u\le u_0\), and \(q\le c_ou^\gamma\), and write \(r=R_P(\pi)\) and \(C_{\mathrm{reg}}=\max\{C_o,1\}>0\).  By \cref{obj:lem:clipped-region-localization},
\[
P_X\!\left(T_\pi(q)\right)
\le
C_{\mathrm{reg}}\,u^\alpha q^{1/\gamma}+\frac{r}{u}.
\]
Both summands are nonnegative, so taking square roots and using \(\sqrt{a+b}\le\sqrt a+\sqrt b\), together with the exponent identity
\[
\sqrt{C_{\mathrm{reg}}\,u^\alpha q^{1/\gamma}}
=
C_{\mathrm{reg}}^{1/2}\,u^{\alpha/2}\,q^{1/(2\gamma)},
\]
gives
\[
P_X\!\left(T_\pi(q)\right)^{1/2}
\le
C_{\mathrm{reg}}^{1/2}\,u^{\alpha/2}q^{1/(2\gamma)}
+
\left(\frac{r}{u}\right)^{1/2}.
\]
% lean: sqrt_add_le_sqrt_add_sqrt % lean: sqrt_margin_overlap_product
Substituting into Step 2 and abbreviating the four nonnegative quantities
\[
J_1=\frac{r_{\mu,n}r_{e,n}}{q},
\qquad
J_2=u^{\alpha/2}q^{1/(2\gamma)},
\qquad
J_3=\left(\frac{r}{u}\right)^{1/2},
\qquad
S=C_{\mathrm{reg}}^{1/2}=\sqrt{\max\{C_o,1\}},
\]
we get \(\left|P[(\pi-\pi^\star_P)b_q]\right|\le 2J_1+2r_{\mu,n}(SJ_2+J_3)\le 4\left(J_1+r_{\mu,n}SJ_2+r_{\mu,n}J_3\right)\), all terms being nonnegative.  Since \(J_1\le(1+S)J_1\), \(Sr_{\mu,n}J_2\le(1+S)r_{\mu,n}J_2\), and \(r_{\mu,n}J_3\le(1+S)r_{\mu,n}J_3\),
\[
4\left(J_1+r_{\mu,n}SJ_2+r_{\mu,n}J_3\right)
\le
4\left(1+\sqrt{\max\{C_o,1\}}\right)\left(J_1+r_{\mu,n}J_2+r_{\mu,n}J_3\right),
\]
that is, with \(C_0=4\{1+\sqrt{\max\{C_o,1\}}\}\) exactly as in the statement,
\[
\left|P\!\left[(\pi-\pi^\star_P)b_q\right]\right|
\le
C_0\left\{
\frac{r_{\mu,n}r_{e,n}}{q}
+
r_{\mu,n}u^{\alpha/2}q^{1/(2\gamma)}
+
r_{\mu,n}\left(\frac{r}{u}\right)^{1/2}
\right\}.
\]
% lean: clipBias_drift_l2_localized_from_region % lean: localized_clipped_drift_bound

\emph{Step 4: the case \(\gamma=0\).}  Assume \(\gamma=0\) and \(q\le\underline p/2\).  \Cref{obj:ass:strict-overlap-endpoint} gives \(\underline p>0\) and \(p_P(X)\ge\underline p\) \(P_X\)-almost surely.  Since \(q\le\underline p/2<\underline p\), the set \(\{p_P\le q\}\) is \(P_X\)-null, hence so is its subset \(T_\pi(q)\):
\[
P_X\!\left(T_\pi(q)\right)=0 .
\]
The second term of the basic drift bound of Step 2 therefore vanishes, leaving
\[
\left|P\!\left[(\pi-\pi^\star_P)b_q\right]\right|
\le
2\,\frac{r_{\mu,n}r_{e,n}}{q}
=
\frac2q\,r_{\mu,n}r_{e,n}
\le
C_1\,r_{\mu,n}r_{e,n},
\qquad
C_1=\max\{1,2/q\},
\]
the last step because \(2/q\le\max\{1,2/q\}\) and \(r_{\mu,n}r_{e,n}\ge0\).  This is the asserted strict-overlap bound. % lean: clipBias_drift_l2_localized_from_region % lean: localized_clipped_drift_bound
\end{proof}
